% Willkommen im Quelltext. :-)
%
% Diese „Bachelor-Arbeit“ wurde mit LaTeX erstellt. Sie kann sowohl
% als Vorlage für Ihre eigene Bachelor-Arbeit verwendet werden als
% auch als Einführung in das Arbeiten mit LaTeX.

% Lizenz: @@@ Tja ... CC-0? CC-BY-SA?

% Im Gegensatz zu What-You-See-Is-What-You-Get- (WYSIWYG-) Systemen
% wie z.B. Textverarbeitungsprogrammen im Office-Bereich, unterscheidet
% LaTeX zwei Versionen des Dokuments: den Quelltext und das compilierte
% Dokument.
%
%  * Der Quelltext ist diejenige Version, die man bearbeitet.
%    Dies geschieht mit einem beliebigen Texteditor, mit dem man
%    auch Computerprogramme schreiben würde. Die LaTeX-Sprache ist
%    eine Markup-Sprache – vergleichbar mit XML oder HTML –, stellt
%    dabei aber mächtigere Werkzeuge zur Verfügung bis hin zu echter
%    Programmierung (Verzweigungen, Schleifen).
%
%  * Durch Compilieren entsteht aus dem Quelltext das druckreife
%    Dokument – typischerweise eine PDF-Datei oder PostScript.
%    Für das vorliegende Dokument geschieht dies durch den Aufruf
%    „pdflatex hauptdatei.tex“ auf der Kommandozeile oder über einen
%    Menüpunkt.
%
% Dieses Dokument enthält zahl- und umfangreiche Kommentare (z.B.
% diesen), die eine kleine Einführung in die Arbeit mit LaTeX geben
% sollen. Kommentare werden mit einem Prozentzeichen eingeleitet und
% gehen jeweils bis zum Zeilenende.
%
% „Befehle“ (genauer: „Macros“) werden in LaTeX typischerweise
% durch einen Backslash eingeleitet und können daher zwischen dem
% eigentlichen Text stehen. Dennoch beginnt ein LaTeX-Dokument nicht
% direkt mit dem Text, sondern es hat einen Vorspann, in dem man
% Pakete einbinden und eigene Macros definieren kann.

% Die \documentclass-Zeile ist in der Regel die erste Zeile eines
% LaTeX-Dokuments. Sie definiert das allgemeine Dokumentenformat:
% Artikel, Bericht (mehrere Kapitel), Buch, Brief usw. Für eine
% Bachelor-Arbeit eignen sich Artikel- und Bericht-Formate. In
% diesem Beispiel verwenden wir das KOMA-Artikel-Format „scrartcl“;
% das entsprechende Bericht-Format heißt „scrreprt“.
%
% Die Optionen [in eckigen Klammern] bedeuten: Verwende A4-Papier,
% setze Gleichungen linksbündig stat zentriert, verwende 12-Punkt-
% statt 11-Punkt-Schrift, trenne die Kopfzeile durch einen
% horizontalen Strich vom Text.
%
\documentclass[a4paper,fleqn,12pt,headsepline]{scrartcl}

% Mit \usepackage bindet man Sammlungen nützlicher Befehle,
% sog. Pakete ein. Das graphicx-Paket stellt z.B. einen Befehl
% \includegraphics zur Verfügung, mit dem man (bei Verwendung von
% PDFLaTeX) Grafiken im PDF-, PNG- oder JPG-Format einbinden kann.
%
\usepackage{graphicx}

% Umlaute und sonstige Sonderzeichen als UTF-8 direkt eingeben
\usepackage[utf8]{inputenc}

% Sprachpaket „reformierte deutsche Rechtschreibung“ einbinden
% (u.a. für Trennungen)
%\usepackage[ngerman]{babel}	% Für Englische Berichte ausqwählen und auskommentieren
\usepackage[english]{babel} 		% Für Englische Berichte ausqwählen und auskommentieren

% Zeichensätze und Zeichensatzkodierung auswählen
%
% Unten sind verschiedene Möglichkeiten aufgeführt, aber nur eine
% wird tatsächlich verwendet. Die anderen sind „auskommentiert“;
% um diese auszuprobieren, genügt es, Kommentarzeichen zu verschieben.
%
\usepackage[T1]{fontenc}
%
%\usepackage{fix-cm} % Standard-Schriftfamilie “Computer Modern” in einer flexibleren Variante verwenden
\usepackage{lmodern} % Schriftfamilie „Latin Modern“ (Erweiterung von „Computer Modern“) verwenden
%\usepackage{droid} % Schriftfamilie „Droid“ verwenden
%\usepackage{times} % „Times New Roman” als Serifenschrift verwenden
%\usepackage{librebaskerville} % „Libre Baskerville” als Serifenschrift verwenden
%\usepackage{libertine} % „Linux Libertine” als Serifenschrift verwenden
%\usepackage{fourier} % „Utopia Regular” als Serifenschrift verwenden
%\usepackage{fouriernc} % „New Century Schoolbook” als Serifenschrift verwenden
%\usepackage{tgschola} % „TeX Gyre Schola” als Serifenschrift verwenden
%\usepackage[scaled=0.95]{helvet} % leicht verkleinerte „Helvetica“ als Sans-Serif-Schrift verwenden
%
% Sans-Serif-Schrift als Standardschrift benutzen
%\renewcommand*\familydefault{\sfdefault}

% Wenn man Anführungszeichen unten und oben nicht von Hand eingeben
% will oder kann (z.B. „Test“), ermöglicht das csquotes-Paket u.a.,
% stattdessen \enquote{Test} zu schreiben.
%
%\usepackage[babel,german=quotes]{csquotes}	% Nur für deutsche berichte einkommentieren
% 

% Das float-Paket stellt erweiterte Möglichkeiten zur Plazierung von
% Abbildungen und sonstiger „gleitender“ Inhalte zur Verfügung.
%
\usepackage{float}

% Bilder mit „Bild“ anstatt mit „Abbildung“ bezeichnen
%\renewcaptionname{ngerman}\figurename{Bild} % für deutsche Berichte einkommentiern 
\renewcaptionname{english}\figurename{Fig.} % für englische Berichte einkommentiern   

% Das picins-Paket stellt Befehle zur Verfügung, Bilder in den Text
% zu integrieren und umfließen zu lassen.
%
%\usepackage{picins}

% Bilder, Tabellen usw. sind in LaTeX normalerweise „gleitend“,
% d.h., sie stehen nicht an einer bestimmten Stelle im Text, sondern
% z.B. oben auf der Seite. Wenn man Bilder usw. nicht gleiten läßt,
% sondern sie fest mit dem Text verbindet, ermöglicht dieses Paket,
% ihre Beschriftung (caption) dennoch im Abbildungsverzeichnis
% erscheinen zu lassen und darauf zu verweisen.
%
% Die hier verwendete KOMA-Version scrartlc des Artikel-Formats
% enthält bereits die in diesem Paket enthaltenen Befehle. Sie
% benötigen es daher nur, wenn Sie das Standard-Artikel-Format
% article verwenden.
% 
%\usepackage{capt-of}

% Paket für mehrseitige Tabellen
%\usepackage{longtable}

% Paket für Seitenrandabstände und Einstellung für Seitenränder
\usepackage{geometry}
\geometry{left=3.5cm, right=2cm, top=2.5cm, bottom=2cm}

% Paket für Kästchen mit Spezialeffekten (z.B. Schatten)
%\usepackage{fancybox}

% bricht lange URLs „schön“ um
\usepackage[hyphens,obeyspaces,spaces]{url}

% Paket für Textfarben
\usepackage{color}

% Zusätzliche mathematische Symbole importieren
\usepackage{amsmath}
\usepackage{amssymb}

% erzeugt Inhaltsverzeichnis mit Querverweisen zu den Kapiteln (PDF Version)
\usepackage[bookmarksnumbered,hyperfootnotes=false]{hyperref} 
\hypersetup{colorlinks,citecolor=black,linkcolor=black,urlcolor=black}

% neue Kopfzeilen mit dem fancyhdr-Paket erzeugen
\usepackage{fancyhdr} % Paket laden
\pagestyle{fancy} % eigener Seitenstil
\fancyhf{} % alle Kopf- und Fußzeilenfelder bereinigen
\fancyhead[L]{\nouppercase{\leftmark}} % Kopfzeile links
\fancyhead[C]{} % zentrierte Kopfzeile
\fancyhead[R]{\thepage} % Kopfzeile rechts
\renewcommand{\headrulewidth}{0.4pt} % obere Trennlinie

% Festlegung der Art der Zitierung auf die Havard-Methode: Abkürzung Autor + Jahr.
% Dieser Befehl gehört zum natbib-Paket.
%
\usepackage{natbib}
\bibliographystyle{apalike_german}

% keine Einrückung am Anfang eines Absatzes,
% stattdessen vertikal etwas Platz lassen
%
\usepackage{parskip}

% Paket zur Erweiterung der Tabellen-Befehle (tabular, array)
%\usepackage{array}

% Dieser Befehl würde den zusätzlichen Zwischenraum abschalten,
% den LaTeX normalerweise nach einem Punkt einfügt.
%\frenchspacing

% Paket für Zeilenabstand einbinden
%\usepackage{setspace}
% Dieser Befehl würde den Zeilenabstand auf „1,5“ einstellen
%\onehalfspacing

% Paket für Listings (z.B. Computerprogramme)
\usepackage{listings}
%
% Einstellungen für C-Quelltexte
%\lstset{numbers=left,
%        numberstyle=\tiny,
%        numbersep=5pt,
%        keywordstyle=\color{black}\bfseries,
%        stringstyle=\ttfamily,
%        showstringspaces=false,
%        basicstyle=\footnotesize,
%        captionpos=b}
%\lstset{language=c}
%
% Einstellungen für LaTeX-Quelltexte
\lstset{basicstyle=\ttfamily,
        language=tex}

% Nomenklatur
%
% Paket einbinden
\usepackage{nomencl}
%
% Befehl „\nomenclature” alternativ als „\abk“ zur Verfügung stellen
\let\abk\nomenclature
%
% Deutsche Überschrift für Nomenklaturverzeichnis
\renewcommand{\nomname}{Nomenklatur} % nur für deutsche Berichte 
%
% Nomenklaturverzeichnis gestalten
% Raum, der den Abkürzungen zur Verfügung steht = 0.2 der Gesamtbreite
\setlength{\nomlabelwidth}{0.2\hsize}
% Zeilenabstände verkleinern
\setlength{\nomitemsep}{-\parsep}
\makenomenclature
% Erstellen von Unterteilungen in der Nomenklatur
\RequirePackage{ifthen}
\renewcommand{\nomgroup}[1]{%
  \ifthenelse{\equal{#1}{A}}{\item[\textbf{Abkürzungen}]\vspace*{1cm}}{%
    \ifthenelse{\equal{#1}{S}}{%
      \item[\textbf{Symbole}]\vspace*{1cm}}{}%
  }%
}

% Benutzte in Überschrift die gleiche Schriftart und Farbe wie im Text
\setkomafont{disposition}{\normalcolor\bfseries}

% Maßeinheit für picture-Grafiken
\setlength{\unitlength}{1cm}

%
% fuer Stichwortverzeichnis
\usepackage{makeidx}
% Stichwortverzeichnis erstellen
\makeindex
\renewcommand{\indexname}{Stichwortverzeichnis} % nur für deutsche Berichte 

%%%%%%%%%%%%%%%
% pgfplots und tikz
\usepackage{tikz}
\usepackage{pgfplots}
\usetikzlibrary{patterns}
\usetikzlibrary{datavisualization}
\usetikzlibrary{arrows}

\usepackage{gensymb}
\usepackage{booktabs}
\usepackage{longtable}
\usepackage{multirow}
\usepackage[toc,page]{appendix} % für englische Berichte


% Bildnummern aus Nummer des Abschnitts und laufender Nummer zusammensetzen
%
% Bei Verwendung eines Bericht-Formats (s.o.) geschieht dies
% automatisch mit der Kapitelnummer (\chapter – eine Ebene oberhalb
% von \section). Für das Artikel-Format machen wir es hier manuell.
%
\renewcommand{\thefigure}{\arabic{section}.\arabic{figure}}
\makeatletter\@addtoreset{figure}{section}\makeatother

% Gleichungsnummern aus Nummer des Abschnitts und laufender Nummer zusammensetzen
\renewcommand{\theequation}{\arabic{section}.\arabic{equation}}
\makeatletter\@addtoreset{equation}{section}\makeatother

% Tabellennummern aus Nummer des Abschnitts und laufender Nummer zusammensetzen
\renewcommand{\thetable}{\arabic{section}.\arabic{table}}
\makeatletter\@addtoreset{table}{section}\makeatother



% auf römische Seitenzahlen umschalten
%\pagenumbering{Roman}

% Mit \input wird eine selbst erstellte .tex-Datei eingebunden.
%
% „trennung.tex“ enthält ein (selbst erstelltes) Verzeichnis mit
% Trennhilfen für häufige Wörter, deren Silbentrennung von der
% Automatik nicht korrekt erkannt wird.
%
%hier müssen alle Wörter rein, welche Latex von sich auch nicht korrekt trennt bzw. bei denen man die genaue Trennung vorgeben möchte
\hyphenation{
Film-pro-du-zen-ten
Lux-em-burg
Soft-ware-bau-steins
zeit-in-ten-siv
wis-sen-schaft-li-chen
ent-spre-chend
}

% Mit \newcommand definieren wir einen neuen Befehl.
% Die 1 in eckigen Klammern bedeutet, daß der Befehl einen Parameter
% (bei Verwendung in geschweiften Klammern zu setzen) akzeptiert.
% Mit #1 nehmen wir auf diesen Parameter bezug.
%
% Der Befehl bewirkt, daß der Parameter in Sans-Serif-Schrift gesetzt wird.
%
% Beispiel für die Verwendung des Befehls: \file{hauptdatei.tex}
%
% Man könnte natürlich auch direkt \textsf{hauptdatei.tex} schreiben.
% Der neue Befehl hat aber den Vorteil, daß er markiert, was wir meinen
% und nicht, wie es aussehen soll. Konkret: Wenn wir uns im Nachhinein
% dafür entscheiden, Dateinamen doch nicht in Sans-Serif-Schrift, sondern
% z.B. kursiv zu setzen, muß dafür nur die Befehlsdefinition geändert werden
% – und nicht jede einzelne Stelle, an der der Befehl verwendet wird.
%
\newcommand{\file}[1]{\textsf{#1}}

% Mit \begin{document} endet der Vorspann, und das eigentliche
% Dokument beginnt.

\begin{document}
%
  % Titelseite
  %
  % Mit \newcommand definieren wir eigene Befehle – in diesem Fall
  % mit dem Zweck, Form und Inhalt zu trennen: In der Datei
  % „deckblatt.tex“ gestalten wir ein Deckblatt zur allgemeinen
  % Verwendung; für den konkreten Titel der Arbeit, Namen des Autors
  % usw. verwenden wir „Variablen“ \thetitle, \theauthor usw., die
  % wir hier vorab setzen.
  %
  \newcommand{\institutname}{  Institut für Mechatronik\\  und Fahrzeugtechnik}
%  \newcommand{\institutname}{  Institut für Informatik\\  und eingebettete Systeme}
  \newcommand{\thetitle}{Optimierung einer Dokumentation\\ % Umbruch an sinnvollen Stellen manuell vornehmen
                         von mechatronischen Beispielprojekten}
  \newcommand{\theauthor}{Max Mustermann}
  \newcommand{\studentnumber}{123\,456\,7890}
  \newcommand{\discipline}{Mechatronik \& Informationstechnologie}
%  \newcommand{\discipline}{Maschinenbau}
%  \newcommand{\discipline}{Elektrotechnik}
%  \newcommand{\discipline}{Mechatronik}
%  \newcommand{\discipline}{Elektromobilität}
  \newcommand{\abschluss}{Bachelor of Engineering}
%  \newcommand{\abschluss}{Master of Engineering}
%  \newcommand{\dokumentname}{Dokumentation des KIS-Projekt-1 im Studiengang}
%  \newcommand{\dokumentname}{Dokumentation des KIS-Projekt-2 im Studiengang}
%  \newcommand{\dokumentname}{Dokumentation des KIS-Projekt-3 im Studiengang}
%  \newcommand{\dokumentname}{Dokumentation der Praxisphase im Studiengang}
  \newcommand{\dokumentname}{Bachelor-Arbeit im Studiengang}
%  \newcommand{\dokumentname}{Master-Arbeit im Studiengang}
  \newcommand{\finaldate}{35.~Mai 2038}
 % Sperrvermerk, falls notwendig 
%
\vspace*{10ex}

\begin{center}
\huge{\textbf{Sperrvermerk}}
\end{center}

\bigskip 


Die vorliegende Bachelor-Thesis mit dem Titel:

\begin{center}
\thetitle
\end{center}

beinhaltet interne und vertrauliche Informationen der Firma:

<Name und Anschrift der Firma>

Die Weitergabe des Inhalts der Arbeit und eventuell beiliegender Zeichnungen und Daten – im Ganzen oder in Teilen – ist bis zum 

<Datum, an dem der Sperrvermerk abläuft> 

grundsätzlich untersagt. Es dürfen keinerlei Kopien oder Abschriften – auch in digitaler Form – gefertigt werden.
Ausnahmen bedürfen der schriftlichen Genehmigung der o. g. Firma.
\bigskip 


Heiligenhaus, den 
%
\underline{\hspace*{2cm}} 
\underline{\hspace*{7cm}}\\
\hspace*{5.5cm}
<Unterschrift der/des Studierenden>

\clearpage
%
 \thispagestyle{empty}

\begin{minipage}{0.55\textwidth}
\institutname{}
%  Institut für Mechatronik\\
%  und Fahrzeugtechnik
\end{minipage}\hfill
%\includegraphics[height=2cm]{bilder/bologo} % Pixelgrafik-Version
\begin{picture}(0,0)(0,1.37) % Vektorgrafik-Version
  \put(-2.7,0.0){\includegraphics[scale=1.06]{bilder/logo-hochschule-bochum-bo.pdf}}
  \put(-6.2,0.87){\includegraphics[scale=0.545]{bilder/logo-hochschule-bochum-text.pdf}}
  \put(-6.2,-0.07){\includegraphics[scale=0.635]{bilder/logo-hochschule-bochum-cvh-text.pdf}}
\end{picture}%

% freier Platz: 1/3 oben 
\vfill

\begin{center}
  % Die äußeren geschweiften Klammern begrenzen die Wirkung von \Huge.
  % Das \par am Ende ist notwendig, damit zum einen der Zeilenabstand
  % des \Huge-Textes angepaßt wird; zum anderen arbeitet \bigskip nicht
  % innerhalb von Textabsätzen, sondern außerhalb.
  {\Huge\textbf{\thetitle}\par}
  \bigskip\bigskip
  von\par
  \bigskip\bigskip
  {\Large\textbf{\theauthor}\par}
  \smallskip
  Matrikelnummer:~\studentnumber\par
  \bigskip\bigskip\bigskip
  \dokumentname{}\\
  \discipline\\
  Software-Entwicklungsprojekt\\
  \abschluss{}
\end{center}

% freier Platz: 2/3 unten
\vfill
\vfill
 
\begin{tabular}{ll}
  \textbf{Eingereicht am:} & \today\\
  \\
  \textbf{1.~Prüfer:} & Prof.~Dr.~rer.nat.~Jörg Frochte \\
  \textbf{2.~Prüfer:} & Benedikt Wildenhain\\

\end{tabular}
 % Inhalt aus separater Datei lesen
%  \thispagestyle{empty}

\begin{minipage}{0.55\textwidth}
\institutname{}
%  Institut für Mechatronik\\
%  und Fahrzeugtechnik
\end{minipage}\hfill
%\includegraphics[height=2cm]{bilder/bologo} % Pixelgrafik-Version
\begin{picture}(0,0)(0,1.37) % Vektorgrafik-Version
  \put(-2.7,0.0){\includegraphics[scale=1.06]{bilder/logo-hochschule-bochum-bo.pdf}}
  \put(-6.2,0.87){\includegraphics[scale=0.545]{bilder/logo-hochschule-bochum-text.pdf}}
  \put(-6.2,-0.07){\includegraphics[scale=0.635]{bilder/logo-hochschule-bochum-cvh-text.pdf}}
\end{picture}%

% freier Platz: 1/3 oben 
\vfill

\begin{center}
  % Die äußeren geschweiften Klammern begrenzen die Wirkung von \Huge.
  % Das \par am Ende ist notwendig, damit zum einen der Zeilenabstand
  % des \Huge-Textes angepaßt wird; zum anderen arbeitet \bigskip nicht
  % innerhalb von Textabsätzen, sondern außerhalb.
  {\Huge\textbf{\thetitle}\par}
  \bigskip\bigskip
  von\par
  \bigskip\bigskip
  {\Large\textbf{\theauthor}\par}
  \smallskip
  Matrikelnummer:~\studentnumber\par
  \bigskip\bigskip\bigskip
  \dokumentname{}\\
  \discipline\\
  mit akademischen Abschluss\\
  \abschluss{}
\end{center}

% freier Platz: 2/3 unten
\vfill
\vfill
 
\begin{tabular}{ll}
  \textbf{Eingereicht am:} & \finaldate\\
  \\
  \textbf{1.~Prüfer:} & Prof.~Dr.-Ing.~Markus Lemmen \\
  \textbf{2.~Prüfer:} & ????
\end{tabular}
 % Inhalt aus separater Datei lesen

  % Abstract
  %
  % Mit \clearpage eröffnen wir eine neue Seite. Wenn sich noch
  % Bilder für den Beginn der nächsten Seite o.ä. im Speicher
  % befinden, soll LaTeX diese jetzt abarbeiten.
  \clearpage
  %
\pagenumbering{Roman}
\setcounter{page}{2}
  % Normalerweise wird bei einem neuen Abschnitt (\section{...}) die
  % Kopfzeile links automatisch auf den Namen dieses Abschnitts gesetzt.
  % Für das Abstract verwenden wir einen „stummen“ Abschnitt
  % (\section*{...}), bei dem dies nicht automatisch geschieht.
  % Mit dem folgenden Befehl setzen wir stattdessen die Kopfzeile
  % „von Hand“ auf „Abstract“.
  \fancyhead[L]{Abstract}
  %
  % Inhalt aus separater Datei lesen
  \section*{Abstract}


Ornithologen bzw. Vogelkundler nutzen das Monitoring von Vögeln für den Vogeltierschutz.

Monitoring ist die dauerhafte Überwachung von bestimmten Systemen.

Im Bereich der Biologie findet sie Anwendung bei Tier- und Pflanzenarten für die Erfassung der Bestände. Es ist durch solche Beobachtungen möglich, Pflanzen und Tiere zu erforschen. Die erfassten Monitoring-Daten helfen Biologen und Umweltschützern den Arten- und Lebensraum nachzuvollziehen und zu schützen .

Diese Arbeit stellt ein KI-gestütztes Tool vor, welches das Monitoring von Vögeln automatisiert unterstützt.

Über ein Mikrofon kann ein Computer mindestens 24 Stunden lang kontinuierlich die Umgebung aufnehmen. Die Software analysiert die Audioaufnahmen basierend auf einem Netzwerk, welches auf die Erkennung von Vogelarten anhand des Vogelgesangs trainiert ist. Anschließend selektiert die Software aus den Audioaufnahmen einzelne Abschnitte und isoliert daraus den erkannten Vogelgesang.

 

%Der Webbrowser?


%https://vogelmonitoring-rlp.de/portfolio/monitoring-wofuer

% Quelle: https://de.wikipedia.org/wiki/Ornithologie




  % Inhaltsverzeichnis auf neuer Seite anzeigen
  \clearpage
%  \fancyhead[L]{Inhaltsverszeichnis}
  \fancyhead[L]{Table of Contents}
  \tableofcontents

  % Abkürzungsverzeichnis
  \clearpage
  \printnomenclature
%  \fancyhead[L]{Nomenklatur} % Kopfzeile links, deutsche Berichte 
  \fancyhead[L]{Nomenclature} % Kopfzeile links, englische Berichte 
  %
  % Im Gegensatz zu anderen Abschnitten erscheint das
  % Abkürzungsverzeichnis nicht automatisch im Inhaltsverzeichnis.
  % Mit dem folgenden Befehl können wir es „von Hand“ ergänzen.
  %\addcontentsline{toc}{section}{Nomenklatur}

  % Abschnitt 1: Einleitung
  \clearpage
   %
  % Wenn wir oben die Kopfzeile links „von Hand“ gesetzt haben,
  % müssen wir sie jetzt auf „Automatik“ zurücksetzen.
  \fancyhead[L]{\nouppercase{\leftmark}}
  %
  % Wenn wir oben auf römische Seitenzahlen umgeschaltet haben,
  % können wir hier auf arabische Seitenzahlen zurückschalten und
  % den Seitenzähler zurücksetzen.
  %\pagenumbering{arabic}
  %
  % den eigentlichen Inhalt aus separater Datei einbinden
  \pagenumbering{arabic}
%
  \section{Einleitende Übersicht%
         \label{sec:Einleitung}}

Thema dieses Projekts ist die Entwicklung einer Software zur praktischen Anwendung in der Erkennung von Vogelarten anhand ihres Gesangs.

Das Projekt soll in der Lage sein, ununterbrochen 24 stundenlang Vogelstimmen in der freien Natur aufzunehmen und zu analysieren. Dieses Verfahren ermöglicht es Ornithologen (aktive Menschen im Bereich Vogelkunde) verschiedene Informationen im Bereich Ökologie und Taxonomie zu erlangen. Das Projekt ist ein Hilfsmittel für spezielle Untersuchungen. Diese Untersuchungen können das Monitoring

\begin{itemize}

\item von den jeweiligen verteilten Populationen

\item des Vogelzugs

\item der versprengten Arten

\end{itemize}

sein.\bigskip

Das Programm soll folgende Kriterien erfüllen:

\begin{itemize}

\item Starten des Programms per Remote-Zugriff

\item Aufnahme und Analyse laufen mindestens 24 Stunden ununterbrochen parallel

\item Aufnahmen, die Vogelgesang beinhalten, für statistische Auswertungen wie Vogelzählung separat speichern und auswerten

\end{itemize}
\bigskip

Die vorliegende Arbeit ist folgendermaßen aufgebaut:

\begin{enumerate}

\item Abschnitt \ref{sec:Konzept} stellt das Konzept dar

\item Abschnitt \ref{sec:versuchsaufbau} skizziert und erklärt den theoretischen Versuchsaufbau

\item Abschnitt \ref{sec:umsetzung} erklärt die Umsetzung (Hardware/Software)

\item Abschnitt \ref{sec:versuchsdurchfuehrung} demonstriert den durchgeführten Versuch

\item Abschnitt ... validiert die Ergebnisse aus den Versuchsdurchführungen %soll es dazu einen extra ABschnitt geben?

\item in Abschnitt \ref{sec:zusammen} sind die Ergebnisse validiert und gibt einen Ausblick auf geschaffene Möglichkeiten durch dieses Projekt sowie mögliche weitere Entwicklungen des Projektes

\end{enumerate}

 


% https://de.wikipedia.org/wiki/Taxonomie

 %https://www.nabu.de/natur-und-landschaft/naturschutz/weltweit/naturschutzprojekte/afribirds/index.html
%https://www.nabu.de/natur-und-landschaft/naturschutz/weltweit/naturschutzprojekte/afribirds/27164.html




  % Abschnitt 2: Grundlagen
  \clearpage
  \section{BirdNET}
- was ist BirdNET?
- wieso gibt es Birdnet
-was kann BirdNET

weiter Informationen unter \cite{kahl2021birdnet}



\section{Audio}

%Quelle: https://wiki.ubuntuusers.de/PulseAudio/


%Quelle: https://wiki.ubuntuusers.de/projectM/

  % Abschnitt 3: Beispiele
  \clearpage
  \input{inhalt/3_latex-beispiele}



  % Abschnitt 4: Zusammenfassung
  \clearpage
  \input{inhalt/4_zusammen}

%  % Abbildungsverzeichnis
%  \clearpage
%  \listoffigures

  % Literaturverzeichnis
  \clearpage
  % manuell in Inhaltsverzeichnis aufnehmen
%  \addcontentsline{toc}{section}{Literatur- und Quellenverzeichnis} % für deutsche Berichte 
  \addcontentsline{toc}{section}{References} % für englische Berichte 
  % Umbenennen
%  \renewcommand\refname{Literatur- und Quellenverzeichnis} % für deutsche Berichte
  % Literaturverzeichnis anzeigen
  \bibliography{hauptdatei}
\clearpage

% 
%. Ggf. Bilderverzeichnis
\listoffigures
\clearpage

%. Ggf. Tabellenverzeichnis
\listoftables
\clearpage

%. Ggf. Abkürzungsverzeichnis
% Steht schon vorne.....

%. Ggf. Stichwortverzeichnis
% Stichwortverzeichnis soll im Inhaltsverzeichnis auftauchen
%\addcontentsline{toc}{section}{Stichwortverzeichnis} % für deutsche Berichte
\addcontentsline{toc}{section}{Index}	% für englische Berichte
% Stichwortverzeichnis endgueltig anzeigen
\printindex
\clearpage


 %%%%%%%%%%%%%%%%%%%%%%%%%%%%%%%%%%%%%%%%%%%%%
  % Anhang
  \clearpage
  \appendix
%  \addcontentsline{toc}{section}{Anhang} % Für deutsche Berichte 
  \addcontentsline{toc}{section}{Appendix}	% für englische Berichte 
  % Formatierung des Anhangs
  % Bildnummern aus Nummer des Abschnitts und laufender Nummer zusammensetzen
  %
  % Bei Verwendung eines Bericht-Formats (s.o.) geschieht dies
  % automatisch mit der Kapitelnummer (\chapter – eine Ebene oberhalb
  % von \section). Für das Artikel-Format machen wir es hier manuell.
  %
  \renewcommand{\thefigure}{\Alph{section}.\arabic{figure}}
  \makeatletter\@addtoreset{figure}{section}\makeatother
  % Gleichungsnummern aus Nummer des Abschnitts und laufender Nummer zusammensetzen
  \renewcommand{\theequation}{\Alph{section}.\arabic{equation}}
  \makeatletter\@addtoreset{equation}{section}\makeatother
  % Tabellennummern aus Nummer des Abschnitts und laufender Nummer zusammensetzen
  \renewcommand{\thetable}{\Alph{section}.\arabic{table}}
  \makeatletter\@addtoreset{table}{section}\makeatother
  %%%%
  % Inhalt des Anhangs
  \input{inhalt/anhang}

  % Eidesstattliche Erklärung
  \clearpage
  \fancyhead[L]{Eidesstattliche Erklärung} % Kopfzeile links
  %\addcontentsline{toc}{section}{Eidesstattliche Erklärung}
  %%%%%%%%%%%%%%%%%%%
\section*{Eidesstattliche Erklärung}
%
Ich versichere, dass ich die Arbeit selbständig verfasst und keinen als die angegebenen Quellen und Hilfsmittel benutzt sowie Zitate kenntlich gemacht habe. 
%
\bigskip 

Die Regelungen der geltenden Prüfungsordnung zu Versäumnis, Rücktritt, Täuschung und Ordnungsverstoß habe ich zur Kenntnis genommen. 
%
\bigskip

Diese Arbeit hat in gleicher oder ähnlicher Form keiner Prüfungsbehörde vorgelegen. 
%
\bigskip 

%Ich versichere, die von mir vorgelegte Arbeit selbstständig verfasst zu haben.
%Alle Stellen, die wörtlich oder sinngemäß aus veröffentlichten oder nicht veröffentlichten 
%Arbeiten anderer entnommen sind, 
%habe ich als entnommen kenntlich gemacht. 
%Sämtliche Quellen und Hilfsmittel, die ich für die Arbeit benutzt habe, sind angegeben. 
%Die Arbeit hat mit gleichem Inhalt bzw. in wesentlichen Teilen noch keiner anderen Prüfungsbehörde vorgelegen.
%

\vspace{3cm}
Heiligenhaus, den 
\underline{\hspace*{3cm}}
\hfill 
\begin{tabular}[t]{@{}l@{}}\hline
\hspace*{2.5cm}Unterschrift\hspace*{2.5cm}
\end{tabular}


\end{document}
